从补条件能得到什么? Agmon, Douglis 和 Nirenberg 的工作 \cite{ADN2} 回答了这一问题. 陈述其结果之前, 先引入一些记号. 对 $M\in\mathbb{N}$, 令
\begin{equation}\label{eq:adn-weighted-sobolev-domain-data-spaces}
\begin{aligned}
X_M&=\prod_{j=1}^kH^{M+t_j}(\Omega),\quad Y_M=\prod_{i=1}^kH^{M-s_i}(\Omega),\\
Z_M&=\prod_{l=1}^mH^{M-r_l-1/2}(\partial\Omega).
\end{aligned}
\end{equation}


现在考虑问题 \eqref{eq:weighted-elliptic-general-system} 及边界条件 \eqref{eq:weighted-elliptic-general-boundary-conditions}. 将方程简写为 $\bm{L}\bm{u}=\bm{f}$ 和 $\bm{B}\bm{u}=\bm{g}$, 并用 $\mathcal{A}$ 表示将 $\bm{u}$ 映到 $(\bm{L}\bm{u},\bm{B}\bm{u}|_{\partial\Omega})$ 的算子. 选取权重时, 现在约定对所有 $i$ 和 $j$ 均有 $s_i\leq0$ 和 $t_j\geq0$; 总可以从 $s_i$ 和 $r_l$ 中减去一个常数, 同时将相同常数加到 $t_j$ 上来做到这一点. 令 $t\prime=\max_jt_j$, $M_1=\max(0,\max_lr_l+1)$. 那么有以下结论.

\begin{Theorem}{Agmon, Douglis, Nirenberg}
\label{thm:agmon-douglis-nirenberg-elliptic-estimate}
设 $M\geq M_1$ 为整数. 假设 $\Omega$ 为 $C^{M+t\prime}$ 类有界区域, $L_{ij}$ 的系数属于 $C^{M-s_i}(\overline\Omega)$, 且 $B_{lj}$ 的系数属于 $C^{M-r_l}(\partial\Omega)$. 此外, 假设在整个 $\overline\Omega$ 上椭圆, 并在 $\partial\Omega$ 的每一点满足补条件. 假设 $\bm{f}\in Y_M$ 且 $\bm{g}\in Z_M$. 那么以下结论成立:\begin{enumerate}\item 每个解 $\bm{u}\in X_{M_1}$ 实际上都属于 $X_M$.\item 存在与 $\bm{u}$, $\bm{f}$ 和 $\bm{g}$ 无关的通用常数 $K$, 使得对每个解 $\bm{u}\in X_M$ 都有
\begin{equation}\label{eq:adn-elliptic-system-apriori-estimate}
\|\bm{u}\|_{X_M}\leq K\left(\|\bm{f}\|_{Y_M}+\|\bm{g}\|_{Z_M}+\sum_{j=1}^k\|u_j\|_{L^2(\Omega)}\right).
\end{equation}
\end{enumerate}若 $\bm{u}$ 为唯一解, 则可省去 \eqref{eq:adn-elliptic-system-apriori-estimate} 中的最后一项.
\end{Theorem}

因而该结果由正则性陈述和先验估计组成. Agmon, Douglis 和 Nirenberg 实际证明的内容比这里陈述的更多; 他们在基于 $L^p$ 的 Sobolev 空间以及 Hölder 空间中都建立了类似结果. 还应注意, 关于 $\Omega$ 和系数的某些光滑性假设可以减弱. 这里不进一步讨论这一点. 定理的证明超出本入门教材的范围. 不过, 一个特殊情形, 即具有 Dirichlet 边界条件的二阶椭圆 PDE, 的证明参见 Sections~\ref{sec:elliptic-interior-regularity} 和 \ref{sec:elliptic-boundary-regularity}. 下面导出一个有趣的推论.

\begin{Corollary}
\label{cor:adn-boundary-operator-semi-fredholm}
作与前一定理相同的全部假设. 再假设对每个 $j$ 都有 $M+t_j>0$. 那么算子 $\mathcal{A}:X_M\to Y_M\times Z_M$ 为半 Fredholm 算子.
\end{Corollary}

\begin{Proof}
由系数的光滑性假设容易得到, $\mathcal{A}$ 确实将 $X_M$ 映到 $Y_M\times Z_M$. 设 $N(\mathcal{A})$ 为 $\mathcal{A}$ 的零空间, $B$ 为 $N(\mathcal{A})$ 与 $(L^2(\Omega))^k$ 中单位球的交集. 由该定理, $B$ 在 $X_M$ 的范数下有界, 因而在 $(L^2(\Omega))^k$ 中预紧. 由于无穷维空间中的单位球绝不可能预紧, 所以 $N(\mathcal{A})$ 必有限维.

接着证明 $\mathcal{A}$ 的值域是闭的. 为此, 假设 $\bm{u}_N$ 解 $\bm{L}\bm{u}_N=\bm{f}_N$, 满足边界条件 $\bm{B}\bm{u}_N=\bm{g}_N$, 并且 $\bm{f}_N$ 和 $\bm{g}_N$ 分别在 $Y_M$ 和 $Z_M$ 中收敛到 $\bm{f}$ 和 $\bm{g}$. 不失一般性, 可假设 $\bm{u}_N$ 在 $(L^2(\Omega))^k$ 中垂直于 $N(\mathcal{A})$. 断言此时 $\bm{u}_N$ 在 $(L^2(\Omega))^k$ 中有界. 反设不然. 抽取子序列后, 可假设 $\|\bm{u}_N\|_2\to\infty$. 令 $\bm{v}_N=\bm{u}_N/\|\bm{u}_N\|_2$. 那么 $\bm{v}_N$ 解问题 $\bm{L}\bm{v}_N=\bm{f}_N/\|\bm{u}_N\|_2$, 并满足边界条件 $\bm{B}\bm{v}_N=\bm{g}_N/\|\bm{u}_N\|_2$. 由 \eqref{eq:adn-elliptic-system-apriori-estimate}, 序列 $\bm{v}_N$ 在 $X_M$ 中有界. 因而它有一个在 $X_M$ 中弱收敛, 从而在 $(L^2(\Omega))^k$ 中强收敛的子序列. 设 $\bm{v}$ 为极限. 那么 $\bm{v}$ 同时属于 $\mathcal{A}$ 的零空间及其正交补, 因而为零. 但这产生矛盾, 因为 $\|\bm{v}\|_2=\lim_{n\to\infty}\|\bm{v}_N\|_2=1$. 由于 $\bm{u}_N$ 在 $(L^2(\Omega))^k$ 中有界, \eqref{eq:adn-elliptic-system-apriori-estimate} 说明它在 $X_M$ 中也有界. 因而抽取子序列后, $\bm{u}_N$ 在 $X_M$ 中弱收敛, 并在 $(L^2(\Omega))^k$ 中强收敛. 再次应用 \eqref{eq:adn-elliptic-system-apriori-estimate}, 得到 $\bm{u}_N$ 实际在 $X_M$ 中强收敛. 极限 $\bm{u}$ 为 $\bm{L}\bm{u}=\bm{f}$ 的解, 满足边界条件 $\bm{B}\bm{u}=\bm{g}$.
\end{Proof}
